\documentclass[10pt,a4paper]{amsart}
\usepackage{amsmath,amssymb,amsthm,graphicx}
\usepackage{verbatim}
\newenvironment{reference}{\comment}{\endcomment}
%\newenvironment{reference}{}{}
\newenvironment{slogan}{\comment}{\endcomment}
\newenvironment{history}{\comment}{\endcomment}
% For commutative diagrams we use Xy-pic
\usepackage[all]{xy}

% We use 2cell for 2-commutative diagrams.
\xyoption{2cell}
% Theorem environments.
%
\theoremstyle{plain}
\newtheorem{theorem}[subsection]{Theorem}
\newtheorem{proposition}[subsection]{Proposition}
\newtheorem{lemma}[subsection]{Lemma}

\theoremstyle{definition}
\newtheorem{definition}[subsection]{Definition}
\newtheorem{example}[subsection]{Example}
\newtheorem{exercise}[subsection]{Exercise}
\newtheorem{situation}[subsection]{Situation}

\theoremstyle{remark}
\newtheorem{remark}[subsection]{Remark}
\newtheorem{remarks}[subsection]{Remarks}

\numberwithin{equation}{subsection}

% Macros
%
\def\lim{\mathop{\mathrm{lim}}\nolimits}
\def\colim{\mathop{\mathrm{colim}}\nolimits}
\def\Spec{\mathop{\mathrm{Spec}}}
\def\Hom{\mathop{\mathrm{Hom}}\nolimits}
\def\Ext{\mathop{\mathrm{Ext}}\nolimits}
\def\SheafHom{\mathop{\mathcal{H}\!\mathit{om}}\nolimits}
\def\SheafExt{\mathop{\mathcal{E}\!\mathit{xt}}\nolimits}
\def\Sch{\mathit{Sch}}
\def\Mor{\mathop{\mathrm{Mor}}\nolimits}
\def\Ob{\mathop{\mathrm{Ob}}\nolimits}
\def\Sh{\mathop{\mathit{Sh}}\nolimits}
\def\NL{\mathop{N\!L}\nolimits}
\def\CH{\mathop{\mathrm{CH}}\nolimits}
\def\proetale{{pro\text{-}\acute{e}tale}}
\def\etale{{\acute{e}tale}}
\def\QCoh{\mathit{QCoh}}
\def\Ker{\mathop{\mathrm{Ker}}}
\def\Im{\mathop{\mathrm{Im}}}
\def\Coker{\mathop{\mathrm{Coker}}}
\def\Coim{\mathop{\mathrm{Coim}}}

% Boxtimes
%
\DeclareMathSymbol{\boxtimes}{\mathbin}{AMSa}{"02}

%
% Macros for moduli stacks/spaces
%
\def\QCohstack{\mathcal{QC}\!\mathit{oh}}
\def\Cohstack{\mathcal{C}\!\mathit{oh}}
\def\Spacesstack{\mathcal{S}\!\mathit{paces}}
\def\Quotfunctor{\mathrm{Quot}}
\def\Hilbfunctor{\mathrm{Hilb}}
\def\Curvesstack{\mathcal{C}\!\mathit{urves}}
\def\Polarizedstack{\mathcal{P}\!\mathit{olarized}}
\def\Complexesstack{\mathcal{C}\!\mathit{omplexes}}
% \Pic is the operator that assigns to X its picard group, usage \Pic(X)
% \Picardstack_{X/B} denotes the Picard stack of X over B
% \Picardfunctor_{X/B} denotes the Picard functor of X over B
\def\Pic{\mathop{\mathrm{Pic}}\nolimits}
\def\Picardstack{\mathcal{P}\!\mathit{ic}}
\def\Picardfunctor{\mathrm{Pic}}
\def\Deformationcategory{\mathcal{D}\!\mathit{ef}}
\usepackage{cleveref}
\crefname{originallemma}{original lemma}{original lemmas}
\Crefname{originallemma}{Original lemma}{Original lemmas}
\makeatletter
\newlabel{lemma-RHom-is-tensor-special}{{lemma-RHom-is-tensor-special}{source}{}{sourcepointer-1}{}}
\newlabel{lemma-RHom-is-tensor-special@cref}{{[originallemma][0][]lemma-RHom-is-tensor-special}{[1][0][]source}{}{}{}}
\makeatother
\crefname{originalreference}{original reference}{original references}
\Crefname{originalreference}{Original reference}{Original references}
\makeatletter
\newlabel{more-algebra-lemma-double-base-change}{{more-algebra-lemma-double-base-change}{source}{}{sourcepointer-2}{}}
\newlabel{more-algebra-lemma-double-base-change@cref}{{[originalreference][0][]more-algebra-lemma-double-base-change}{[1][0][]source}{}{}{}}
\makeatother
\crefname{originalreference}{original reference}{original references}
\Crefname{originalreference}{Original reference}{Original references}
\makeatletter
\newlabel{more-algebra-lemma-flat-finite-presentation-perfect}{{more-algebra-lemma-flat-finite-presentation-perfect}{source}{}{sourcepointer-3}{}}
\newlabel{more-algebra-lemma-flat-finite-presentation-perfect@cref}{{[originalreference][0][]more-algebra-lemma-flat-finite-presentation-perfect}{[1][0][]source}{}{}{}}
\makeatother
\crefname{originalreference}{original reference}{original references}
\Crefname{originalreference}{Original reference}{Original references}
\makeatletter
\newlabel{more-algebra-lemma-perfect-ring-map}{{more-algebra-lemma-perfect-ring-map}{source}{}{sourcepointer-4}{}}
\newlabel{more-algebra-lemma-perfect-ring-map@cref}{{[originalreference][0][]more-algebra-lemma-perfect-ring-map}{[1][0][]source}{}{}{}}
\makeatother
\title{The upper shriek of a perfect homomorphism is tensoring with its value on the structure sheaf}
\author{The Stacks Project Authors}
\date{stacks-project:dualizing.tex:local-snapshot-20261012; GNU FDL 1.2 or later}
\begin{document}
\maketitle
\noindent\textit{Editorial note.} All delivered mathematical text is the author's own, copied from the pinned Stacks Project chapter source sources/stacks-dualizing/pinned-original-source.tex; no statement or proof step was rewritten, reordered or removed.\par\medskip
\noindent\textit{Editorial note.} Preamble provenance: the Stacks Project keeps one single shared preamble for every chapter; this card takes it directly from the staged sources/stacks-dualizing/preamble.tex as three verbatim spans (lines 8-12, 14-18 and 154-217), so no derived file was written and nothing was added or redefined. The wrapper class and the base packages are supplied by the assembler.\par\medskip
\noindent\textit{Editorial note.} Bibliography provenance: neither the counted statement nor its proof contains a citation (this lane rescans the whole span before writing the card), so the card carries no bibliography and no entry was reconstructed.\par\medskip
\noindent\textit{Editorial note.} Selection provenance: this card was opened under the file-based control contract (CONTROL.md), with the mechanical pre-card gate run on its statement and proof interval before the card was written.\par\medskip
\noindent\textit{Allowed standard prerequisites (editorial).}\begin{itemize}
\item The perfectness of the target as a module over the chosen polynomial ring in the perfect case. Invoked inside the counted proof (source lines 4901-4903). Supplies the hypothesis that makes the tensor description of the internal hom available.
\item The description of the internal hom out of a perfect module as tensoring with the internal hom of the module. Invoked inside the counted proof (source lines 4905-4906). Rewrites the internal hom defining the shriek as a tensor product.
\item The double base change identity for the derived tensor product along the two presentations. Invoked inside the counted proof (source lines 4905-4907). Reassembles the tensor factors into tensoring with the shriek of the structure sheaf.
\end{itemize}\medskip
\ifdefined\hypertarget\hypertarget{sourcepointer-1}{}\fi
\noindent\textit{Source reference.} Original lemma lemma-RHom-is-tensor-special: sources/stacks-dualizing/pinned-original-source.tex, source lines 2582--2582. Navigation pointer only: the counted text refers to the Stacks item lemma-RHom-is-tensor-special, which is an existing labeled result of this same chapter and is not part of this card. The pointer carries the target's own label as its number and its definition line as its span; its content is not reproduced here.\par\medskip
\ifdefined\hypertarget\hypertarget{sourcepointer-2}{}\fi
\noindent\textit{Source reference.} Original reference more-algebra-lemma-double-base-change: sources/stacks-dualizing/pinned-original-source.tex, source lines 4695--4695. Navigation pointer only: the counted text refers to the Stacks item more-algebra-lemma-double-base-change, which belongs to another chapter of the project and is not part of this card. The pointer carries the target's own label as its number because this chapter's staged snapshot does not contain the other chapter's tag table; its content is not reproduced here.\par\medskip
\ifdefined\hypertarget\hypertarget{sourcepointer-3}{}\fi
\noindent\textit{Source reference.} Original reference more-algebra-lemma-flat-finite-presentation-perfect: sources/stacks-dualizing/pinned-original-source.tex, source lines 4771--4771. Navigation pointer only: the counted text refers to the Stacks item more-algebra-lemma-flat-finite-presentation-perfect, which belongs to another chapter of the project and is not part of this card. The pointer carries the target's own label as its number because this chapter's staged snapshot does not contain the other chapter's tag table; its content is not reproduced here.\par\medskip
\ifdefined\hypertarget\hypertarget{sourcepointer-4}{}\fi
\noindent\textit{Source reference.} Original reference more-algebra-lemma-perfect-ring-map: sources/stacks-dualizing/pinned-original-source.tex, source lines 4650--4650. Navigation pointer only: the counted text refers to the Stacks item more-algebra-lemma-perfect-ring-map, which belongs to another chapter of the project and is not part of this card. The pointer carries the target's own label as its number because this chapter's staged snapshot does not contain the other chapter's tag table; its content is not reproduced here.\par\medskip
\begin{lemma}
\label{lemma-upper-shriek-is-tensor-functor}
Let $\varphi : R \to A$ be a perfect homomorphism of Noetherian rings
(for example $\varphi$ is flat of finite type).
Then $\varphi^!(K) = K \otimes_R^\mathbf{L} \varphi^!(R)$
for $K \in D(R)$.
\end{lemma}

\begin{proof}
(The parenthetical statement follows from
More on Algebra, Lemma
\ref{more-algebra-lemma-flat-finite-presentation-perfect}.)
We can choose a factorization $R \to P \to A$ where $P$ is a polynomial
ring in $n$ variables over $R$ and then $A$ is a perfect $P$-module, see
More on Algebra, Lemma \ref{more-algebra-lemma-perfect-ring-map}.
Recall that $\varphi^!(K) = R\Hom(A, K \otimes_R^\mathbf{L} P[n])$.
Thus the result follows from
Lemma \ref{lemma-RHom-is-tensor-special}
and More on Algebra, Lemma \ref{more-algebra-lemma-double-base-change}.
\end{proof}

\end{document}
